%% =====================================================================
%%  B1-T-16-01 -- what B1 contributes to "Concession and Retreat"
%%  -------------------------------------------------------------------
%%  LAYER A: APPEND-ONLY. Written once at the entry's write, then left
%%  alone. If the source has more to say later, add a bullet -- do not
%%  rewrite. Material found ONLY here goes in the `unique` slot with
%%  @unique set to yes, which makes it undeletable (rule R10).
%% =====================================================================
%% @kind: source
%% @id: B1-T-16-01
%% @book: B1
%% @topic: T-16-01
%% @locator: Tactic 10 (ch. XI), pp. 117--132 (additive)
%% @status: distilled
%% @unique: no
%% @integrated: yes
%% @updated: 2026-09-19

\dossier{B1}{T-16-01}{\bookshort{B1}}{Tactic 10 (ch. XI), pp. 117--132 (additive)}

\begin{distilled}
  \item B1 scopes the concession: the strength to let a useless argument go --- especially when you are right --- is named a mark of maturity, and most arguments are useless by this test.
  \item Concede only where no vital interest (mental, physical, financial) is engaged; the sleeping-dogs rule decides which fights should never have begun.
  \item The chapter's arithmetic: winning the argument that changes nothing is a net loss of composure, time and standing.
\end{distilled}
